\section{Hipóteses de mecanismos para a abordagem}

A questão de projeto é quais intervenções permitem satisfazer os critérios da Seção~\ref{sec:protocolo}. O espaço de investigação inclui composição e roteamento de modelos, engenharia de contexto, adaptação do processo e verificação. As revisões de Wang et al. e Ferrag et al. situam planejamento, contexto, ferramentas e avaliação como componentes dos sistemas de agentes, sem estabelecer uma configuração universalmente superior \cite{Wang2025AgenticProgramming,Ferrag2026AutonomousAgents}. Aqui, esses componentes serão examinados por seus efeitos sobre resolução e recursos, mantendo os pesos dos modelos fixos e sem pressupor a necessidade de uma arquitetura multiagente.

\subsection{Atividades, modelos e controle do processo}

Uma atividade, como localizar código ou validar um patch, não exige necessariamente um agente autônomo ou um modelo exclusivo. A composição fixa atribui modelos a atividades previamente definidas; o roteamento de modelos escolhe qual modelo acionar a partir dos sinais disponíveis. O roteamento de processo decide quais etapas executar, enquanto a seleção de contexto determina a informação fornecida a cada etapa. Essas decisões podem interagir, mas suas contribuições não serão confundidas com o número de agentes.

FlowGen organiza agentes segundo atividades de desenvolvimento e compara processos inspirados em Waterfall, TDD e Scrum. Em HumanEval e MBPP, a remoção de testing prejudicou fortemente a correção funcional, enquanto os efeitos de outras atividades foram menores ou inconsistentes para essa medida; design e revisão também afetaram indicadores de qualidade estática \cite{Lin2024FlowGen}. Esses resultados sustentam avaliar atividades separadamente, mas sua transferência de problemas predominantemente algorítmicos para issues em repositórios permanece limitada. Já Agentless utiliza um processo fixo de localização, reparo e validação, sem delegar ao modelo a escolha livre da próxima etapa. Sua configuração gera múltiplos patches e testes candidatos, mostrando que simplicidade de controle não implica uma única chamada ou baixo consumo por definição \cite{Xia2024Agentless}.

A adequação do processo depende também do modelo. No SWE-Gym, os modelos-base Qwen de 7B e 32B, sem o ajuste fino estudado, resolvem respectivamente 1,0\% e 3,0\% do SWE-bench Lite com OpenHands, contra 7,0\% e 19,0\% com MoatlessTools, cujo workflow restringe o espaço de ações \cite{Pan2025SWEGym}. O contraste envolve scaffolds completos e não isola a restrição de autonomia. Ele motiva investigar se um processo mais adequado permite aproveitar melhor modelos de menor custo. O mesmo trabalho estuda treinamento de agentes e verificadores; disponibilizar instruções procedurais ou ferramentas constitui uma intervenção em inferência e não será apresentado como substituto experimentalmente equivalente ao ajuste fino.

\subsection{Composição de modelos por atividade}

Jiang et al. distinguem DirectSolve, que gera uma solução a partir de um contexto amplo, de SelectSolve, que separa localização de arquivos e reparo. No SWE-bench Verified, DirectSolve com Gemini 1.5 Pro alcança 38,0\% de resolução, e SelectSolve com esse mesmo modelo nas duas etapas alcança 39,2\%. Mantendo Gemini na seleção e empregando Claude 3.7 Sonnet no reparo, SelectSolve alcança 48,6\% \cite{Jiang2025LCLM}. O contraste motiva a especialização por atividade, mas não demonstra vantagem econômica nem a eficácia de um roteador dinâmico. Os resultados principais incluem amostragem de oito patches e seleção de uma entrega; portanto, não representam a acurácia de uma única geração sem processamento adicional.

A vantagem da separação não é uniforme. Em uma ablação com 100 tarefas, DirectSolve e SelectSolve obtêm, respectivamente, 51\% e 49\% com Gemini 2.5 Pro, enquanto a separação apresenta estimativas maiores para os demais modelos examinados. O estudo também reporta custos agregados superiores para seus métodos de contexto longo, sem fornecer uma comparação detalhada de custo e latência entre todas as composições \cite{Jiang2025LCLM}. Assim, a escolha do modelo e a organização do contexto serão investigadas conjuntamente, sem pressupor que decompor atividades ou combinar provedores seja necessário.

A hipótese inicial de distribuição de papéis combina execução predominantemente por um modelo de menor custo com planejamento, revisão ou diagnóstico por um modelo mais capaz. Uma alternativa inverte essa distribuição: um modelo de menor custo prepara o contexto e o mais capaz implementa a solução. Uma variante empregaria um orquestrador para acompanhar dependências, resultados intermediários, orçamento e encerramento. Essas hipóteses permanecem abertas: especialização e assistência podem melhorar a execução, mas também introduzir custo e falhas de transferência de contexto. Gatilhos de escalonamento, conteúdo da assistência e limites de intervenção são decisões do mecanismo a investigar.

Antes de justificar essa assistência, é necessário compará-la com uma alternativa mais simples: permitir que o modelo barato tente resolver a tarefa novamente. Essa comparação envolve três questões: se o modelo consegue gerar uma solução correta, se o sistema consegue reconhecê-la e se os recursos necessários compensam o benefício obtido.

\subsection{Gerar uma solução correta}

Repetir a geração produz várias respostas candidatas para a mesma tarefa, por exemplo, diferentes propostas de alteração do código. Algumas podem estar corretas e outras não. Em um caso idealizado, com tentativas independentes e probabilidade fixa $p>0$ de acerto em cada uma, a probabilidade de haver ao menos uma solução correta entre as $k$ respostas geradas é $1-(1-p)^k$. Essa probabilidade converge a um quando o número de tentativas cresce indefinidamente. A conclusão diz respeito à existência de uma solução correta entre as respostas candidatas: a expressão não informa qual delas está correta nem como escolhê-la para entrega. Se a probabilidade de acerto for zero, repetir a mesma estratégia de geração não resolve a tarefa. Independência e probabilidade fixa são hipóteses desse caso idealizado, não propriedades asseguradas de um agente que modifica o contexto e recebe orientações entre tentativas.

\subsection{Reconhecer a solução correta}

Para transformar as respostas candidatas em uma entrega, o sistema precisa de um critério de seleção. Testes podem fornecer esse critério, mas uma solução incorreta pode ser aprovada quando o erro não é coberto pelos testes. Essa aprovação indevida constitui um falso positivo. Stroebl et al. observaram, em HumanEval e MBPP com testes de cobertura limitada, que modelos com menor acurácia em uma única resposta também produziram mais falsos positivos. Sob os verificadores estudados, repetir a geração até obter aprovação apresentou um limite de acurácia mesmo com orçamento infinito \cite{Stroebl2026Resampling}. Assim, encontrar uma resposta que passa pelo verificador não equivale necessariamente a encontrar uma resposta correta.

Um caso idealizado explicita essa diferença. Mantendo a probabilidade de acerto $p$, suponha que os pares de geração e julgamento sejam independentes entre tentativas, com probabilidades fixas. Sejam $t$ a probabilidade de o verificador aceitar uma solução correta e $f$ a probabilidade de aceitar uma incorreta. Ao entregar a primeira resposta aprovada, a probabilidade de essa entrega estar correta, no limite de tentativas, é
\begin{equation}
\frac{pt}{pt+(1-p)f},
\end{equation}
quando o denominador é positivo. O numerador representa a probabilidade de gerar uma solução correta e aceitá-la; o denominador inclui também as aceitações de soluções incorretas. Para $0<p<1$ e $f>0$, o limite é inferior a um. Se o verificador não aceitar soluções incorretas ($f=0$) e houver chance positiva de gerar e aceitar uma correta ($pt>0$), o limite será um. Portanto, a confiabilidade da seleção depende tanto do gerador quanto do verificador.

Se é necessário reconhecer erros para decidir quando tentar novamente, uma possibilidade é pedir ao próprio modelo que revise sua resposta. Entretanto, Huang et al. encontraram dificuldades na autocorreção de raciocínio quando a revisão depende apenas do modelo, sem feedback externo; em alguns casos, as revisões pioraram o desempenho \cite{Huang2024SelfCorrect}. Esses resultados não estabelecem impossibilidade de autocorreção em programação: executar o código e observar falhas em testes fornece evidências externas que não estão presentes naquele cenário. A capacidade de produzir uma resposta não deve, portanto, ser tomada como garantia de conseguir avaliar e corrigir essa mesma resposta.

Outra possibilidade é usar um modelo mais capaz como juiz, isto é, encarregá-lo de avaliar a resposta produzida pelo modelo barato. Um juiz que aceite menos soluções incorretas pode melhorar a seleção, mas maior capacidade de geração não assegura julgamento infalível. Também há possíveis dependências entre produção e avaliação: em sumarização, Panickssery et al. identificaram reconhecimento e preferência pelas próprias gerações \cite{Panickssery2024SelfPreference}. A transferência desse resultado para código permanece uma questão empírica, e usar modelos diferentes não assegura que seus erros sejam independentes.

O acesso a evidências pode ser tão relevante quanto a escolha do revisor. Agent-as-a-Judge emprega ferramentas e inspeção dos artefatos para verificar requisitos de desenvolvimento de software \cite{Zhuge2024AgentJudge}. Em busca, sistemas de dados e interfaces gráficas, AJ-Bench encontra ganhos com avaliação baseada em ambiente, mas também falhas na aquisição de evidências e na interpretação dos resultados \cite{Shi2026AJBench}. Esses trabalhos motivam revisão apoiada em testes e artefatos disponíveis, sem presumir que cada etapa adicional compense seu custo ou que seus resultados se transfiram diretamente para resolução de issues.

\subsection{Decidir se a assistência compensa}

Mesmo quando repetição e revisão aumentam a chance de uma entrega correta, cada tentativa, teste e chamada ao juiz consome recursos. Orçamento ilimitado é uma condição teórica; a questão operacional é quanto benefício pode ser obtido dentro dos limites de custo e latência. No estudo de Stroebl et al., o número de tentativas de maior utilidade dependeu da penalização atribuída aos falsos positivos \cite{Stroebl2026Resampling}. Não decorre disso um número universal de repetições, nem que poucas tentativas sejam economicamente desvantajosas. É necessário comparar seu consumo e sua resolução com a execução direta pelo modelo forte e com outras formas de assistência.

Essa assistência pode ir além do julgamento. Um modelo forte pode diagnosticar a falha, orientar uma nova tentativa ou modificar o código. Nesses casos, ele intervém na produção das soluções, alterando as condições do modelo de repetição com probabilidade fixa. O mecanismo deverá distinguir nova geração, revisão de uma resposta existente, reparo e escalonamento, entendido como transferência de parte do trabalho ao modelo forte. A ajuda pode melhorar a solução, mas também consumir mais recursos ou introduzir erros ao transferir informações entre modelos.

Após uma revisão, continuar reparando tampouco garante melhoria. Uma alteração pode corrigir uma falha ou tornar incorreta uma solução antes correta. VRR-Stop distingue essas possibilidades dos erros do verificador e propõe decidir pela continuação com base na melhoria esperada da próxima rodada \cite{Wu2026VRRStop}. Seus resultados em raciocínio, incluindo cenários de estresse com perturbações controladas, fundamentam investigar regras de parada, sem estabelecer ganhos equivalentes em programação. Aqui, a hipótese é que os sinais disponíveis durante a execução permitam escolher entre repetir, reparar, escalar ou encerrar, obtendo benefício suficiente para compensar o gasto adicional e o risco de piorar a solução.

Portanto, ampliar o número de tentativas do modelo barato não assegura, por si só, que o sistema entregue uma solução correta ou alcance a taxa de resolução do modelo forte. O benefício depende também da capacidade de reconhecer erros e acertos e dos recursos necessários para isso. A assistência de outro modelo constitui uma hipótese para melhorar essa relação, a ser comparada com repetição simples e execução direta pelo modelo forte. Sua adoção só se justificará se a combinação de resolução, custo e latência atender aos critérios do protocolo; uma melhoria de resolução isolada não demonstrará vantagem econômica.

\subsection{Aquisição, seleção e distribuição de contexto}

Engenharia de contexto compreende, neste estudo, a aquisição, seleção, representação, distribuição e atualização da informação disponível durante a execução. Seu efeito pode decorrer tanto de incluir evidências ausentes quanto de reduzir interferência ou consultas redundantes. Em uma análise controlada de DirectSolve, mantendo os arquivos-alvo no início do prompt, ampliar o contexto de 100 mil para um milhão de tokens reduziu a medida reportada de pass@1 de 39\% para 31\% \cite{Jiang2025LCLM}. Como a informação necessária já estava presente, o resultado evidencia uma limitação no aproveitamento de contexto adicional, mas não demonstra que um seletor operacional consiga identificar previamente esses arquivos.

A seleção também pode preservar hipóteses alternativas. Em Agentless, gerar 40 patches a partir de localizações reunidas em um único contexto resolveu 85 issues; distribuir o mesmo número de candidatos entre quatro conjuntos de localização resolveu 96. Os custos diferiram e as configurações não isolam uma causa única, mas o contraste motiva investigar contextos alternativos em vez de agregar toda informação recuperada \cite{Xia2024Agentless}. Cobertura e concisão serão tratadas conjuntamente: excluir informação necessária pode reduzir tokens e, ao mesmo tempo, comprometer a solução ou provocar novas consultas.

Conhecimento procedural adicional apresenta limitações semelhantes. Nos resultados preliminares do SWE-Skills-Bench, com Claude Code e Haiku 4.5, 39 de 49 skills não alteraram a taxa de aprovação; 24 dessas já alcançavam 100\% sem skill. Algumas reduziram tokens sem alterar a aprovação, enquanto outras aumentaram o consumo ou prejudicaram o resultado. A média passou de 89,8\% para 91,0\%, com efeitos heterogêneos entre skills \cite{Han2026SWESkills}. O estudo disponibiliza uma skill previamente selecionada por tarefa e deixa seleção dinâmica e composição de várias skills como questões futuras. A hipótese aqui é que selecionar orientação compatível com tarefa, versão e contexto possa melhorar a relação entre resolução e recursos; sua eficácia não será inferida apenas da pertinência temática da skill.

A consulta a estruturas do código oferece outro mecanismo de aquisição. O icat-agent emprega ferramentas de inspeção de símbolos e cadeias de chamadas, mas seus resultados não isolam o benefício dessas ferramentas \cite{Chen2026IcatAgent}. CodeGraph será considerado como candidato de implementação para uma variante de recuperação estrutural, sujeito à verificação de suporte e à comparação com a busca da configuração de referência. A avaliação deverá distinguir custo de indexação, atualização e consulta. A integração de ferramentas via MCP, se adotada, será uma decisão de interface; seu uso não constituirá evidência de melhoria na seleção ou no raciocínio.

\subsection{Ativação de etapas e verificação intermediária}

O icat-agent adapta o processo conforme uma avaliação da descrição da issue: ativa um Explorer quando faltam informações e, nos demais casos, inicia diretamente edição e validação em paralelo. Dispensar o Explorer não elimina a inspeção do repositório pelo editor. Na ablação com 60 issues bem especificadas por benchmark, forçar exploração elevou o custo e não aumentou a resolução agregada. Os subconjuntos continham 30 casos originalmente resolvidos e 30 não resolvidos, de modo que os efeitos não estimam diretamente a vantagem da política no benchmark completo \cite{Chen2026IcatAgent}. Essa evidência motiva uma hipótese restrita: ativar uma etapa dedicada de exploração somente quando seus benefícios esperados compensarem o trabalho adicional, comparando essa política com processos fixos.

O mesmo sistema separa contextos de edição e validação e restringe a comunicação a resultados estruturados, evitando compartilhar diretamente o código dos testes com o editor e o histórico de raciocínio do editor com o validador \cite{Chen2026IcatAgent}. Essa separação pode reduzir dependências indesejadas, mas não garante independência estatística dos erros. A hipótese de que o isolamento reduza aceitações indevidas será examinada por contrastes que explicitem a informação compartilhada, preservando o avaliador reservado como medida externa.

Verificar premissas antes de implementar constitui uma intervenção distinta. Spec Kit Agents introduz consultas ao repositório e verificações de artefatos intermediários. Na comparação entre os workflows Full e Full-Augmented, a nota média do juiz LLM passou de 3,51 para 3,66, enquanto o tempo reportado para pares concluídos passou de 24,0 para 37,2 minutos. A pequena avaliação humana, com seis tarefas, registrou mais preferências pelo Full do que pelo Full-Augmented, além de uma maioria de empates \cite{Taghavi2026SpecKitAgents}. Portanto, grounding intermediário motiva investigação, sem demonstrar benefício uniforme entre avaliadores ou justificar sua execução em toda tarefa.

A taxonomia de Macedo distingue especificação, contexto, papéis, execução, validação e portabilidade em frameworks de desenvolvimento; sua comparação documental, pontuada por um único avaliador, não mede superioridade experimental de qualidade ou produtividade \cite{Macedo2026PromptProcess}. Para esta pesquisa, artefatos intermediários poderão organizar decisões, mas suas afirmações precisarão ser confrontadas com requisitos, dependências e evidências de execução. A existência de um plano coerente não será tratada como confirmação de sua compatibilidade com o repositório.

\subsection{Hipóteses e seleção da abordagem}

As evidências motivam investigar três relações: se a composição de modelos por atividade supera alternativas de modelo único; se a adequação e distribuição do contexto permitem reduzir recursos preservando resolução; e se a ativação seletiva de etapas supera processos fixos. São hipóteses sobre a abordagem a desenvolver, e não resultados estabelecidos pelos estudos anteriores. Também permanece aberta a hipótese de que melhorar o contexto antes de escalar o modelo resolva parte das falhas com menor custo. Sua avaliação dependerá de sinais disponíveis durante a execução e incluirá o gasto de obter esses sinais.

A paralelização permanece uma hipótese para reduzir latência em subtarefas sem dependências conflitantes. Integração, verificação e coordenação consumirão recursos e poderão anular essa vantagem. Revisores baseados em modelo também serão investigados como mecanismos internos, mantendo a avaliação externa como medida de resolução. O benefício de revisar dependerá tanto das falhas detectadas quanto do gasto adicional e de possíveis regressões.

Permanece em aberto se a contribuição relativa do scaffold, a estrutura de ferramentas e controle que organiza a atuação do modelo, diminui com o aumento da capacidade do modelo. O estudo em GAIA não encontrou redução monotônica no nível mais difícil e observou diferenças entre famílias; seu domínio não permite antecipar a direção dessa interação nas tarefas de código aqui consideradas \cite{Starace2026ScaffoldGAIA}. Essa hipótese será examinada apenas na medida em que os modelos e as ablações implementadas permitirem contrastes adequados.

A seleção dos mecanismos deverá justificar cada componente por uma hipótese de efeito sobre resolução, custo ou tempo e permitir sua remoção em ablações pertinentes. A arquitetura resultante será uma intervenção submetida ao protocolo, e sua eficiência permanecerá uma hipótese até a avaliação reservada.
